\chapter{The Entropic-BPI algorithm}
\label{chap:algorithm}

This chapter defines the algorithm of the paper (Section~4, Algorithm~1, and the equations of
Appendix~B: (bonus positive), (optimism\_equation\_positive), (width\_certificate\_positive)
and their negative versions) as an identification algorithm in the episode environment, and
the auxiliary ``ring'' quantities of its analysis. The algorithm knows the reward function
$r$ and the parameters $\beta \ne 0$, $\delta \in (0, 1)$, $\eps > 0$, $s_1$; the transition
kernels are unknown to it.

\textbf{Standing setting and abbreviations.} $r$ is a reward function with values in $[0, 1]$
(Definition~\ref{def:reward_fn}), $\beta \ne 0$, $\delta \in (0, 1)$, $\eps > 0$ and $s_1 \in \Ss$.
All the quantities of Sections~\ref{sec:algorithm_backups}--\ref{sec:algorithm_alg} are
computed after $t$ episodes from the history $\hist_t$ of the first $t$ episodes: they only
depend on the counts $n_h^t(s, a)$ (Definition~\ref{def:counts}), the empirical transitions
$\phat_h^t(\cdot \mid s, a)$ (Definition~\ref{def:emp_trans}), the rates $\alpha, \alpha^*$ of
Definition~\ref{def:rates} and $r$; when $t$, $h$ and $(s, a)$ are clear we abbreviate
$n := n_h^t(s, a)$, $\phat := \phat_h^t(\cdot \mid s, a)$ and write $\phat g$ for
$(\phat_h^t g)(s, a)$. The ring quantities of Section~\ref{sec:algorithm_ring} also use the true
kernels $p_h$ and exist only for the analysis.

\textbf{Lean remark.} Every quantity of the chapter is defined as a function of a history
\texttt{hist : Hist Unit (Policy S A H) (Traj S H) t} (file \texttt{Essakine2026Tight/EV2026/Algorithm.lean},
namespace \texttt{Essakine2026Tight}) and evaluated on a run $(X, Y)$ through
\texttt{histAt X Y t ω} (\texttt{Run.lean}: \texttt{UtildeAt}, \texttt{UlowerAt},
\texttt{ZtildeAt}, \texttt{ZlowerAt}, \texttt{bonusAt}, \texttt{greedyAt}, \texttt{certAt},
\texttt{ringZAt}, \texttt{ringUAt}). The backward recursions are indexed by the number
\texttt{j : ℕ} of steps to the end: \texttt{j = 0} is the terminal step $H + 1$ and the paper's
step $h$ has \texttt{j = H + 1 - h}; the Lean step of a pair is \texttt{stepOf H k : Fin H}, the
step with \texttt{k = j - 1 = H - h} steps after it, and the values at the step \texttt{j} are
computed from the values at \texttt{j - 1} by one application of the one-step functions
\texttt{bonus}, \texttt{stepU}, \texttt{ringUAux} (which take the step data
\texttt{k n r phat Zt Zl} as arguments and are independent of histories).

\section{Bonus and backups}
\label{sec:algorithm_backups}

\begin{definition}[Exploration bonus]\label{def:bonus}
\lean{Essakine2026Tight.bonus, Essakine2026Tight.bonusAt}\leanok
\uses{def:counts, def:emp_trans, def:rates, lem:emp_trans_simplex}
Let $t \ge 0$, $h \le H$ and $(s, a)$ with $n = n_h^t(s, a) \ge 1$, and let $\Zt^t_{h+1}, \Zl^t_{h+1} : \Ss \to \R$
be the optimistic and pessimistic values of the next step (Definition~\ref{def:backups}). The
\emph{bonus} is, for $\beta > 0$ (eq.~(bonus positive) of the paper),
\[
  b_h^t(s, a) := 2\sqrt2\, \sqrt{\Var_{\phat}\big(\Zt^t_{h+1}\big)\, \frac{\alpha^*(n)}{n}}
  + 5\, e^{\beta(H - h)}\, \frac{\alpha(n)}{n}
  + 4 H\, e^{\beta(H - h)}\, \frac{\alpha(n)}{n} ,
\]
and for $\beta < 0$ (eq.~(bonus\_negative)),
\[
  b_h^t(s, a) := 2\sqrt2\, \sqrt{\Var_{\phat}\big(\Zl^t_{h+1}\big)\, \frac{\alpha^*(n)}{n}}
  + 5\, \big(1 - e^{\beta(H - h)}\big)\, \frac{\alpha(n)}{n}
  + 4 H\, \big(1 - e^{\beta(H - h)}\big)\, \frac{\alpha(n)}{n} ,
\]
where $\Var_{\phat}(g) = \Var_{\phat_h^t}(g)(s, a)$ is the variance under the empirical
transitions (Definition~\ref{def:emp_trans}). The bonus is not defined (and never used) for a
pair with $n_h^t(s, a) = 0$. The third term carries the KL rate $\alpha$ (as in the statement of
the paper's Lemma~7), not the Bernstein rate $\alpha^*$ of the paper's equations
(bonus positive) and (bonus\_negative): with $\alpha^*$, which can be as small as $\alpha/S$,
the certificate lemmas (Lemmas~\ref{lem:certificate_pos} and~\ref{lem:certificate_neg}) fail
(Remark~\ref{rem:pos_constants_bonus}).
\end{definition}

\textbf{Lean remark.} \texttt{bonus nA H β δ k n phat Zt Zl : ℝ} with \texttt{k = H - h} the
number of steps after the step (so $e^{\beta(H - h)}$ is \texttt{Real.exp (β * k)}),
\texttt{nA = Fintype.card A}, the rates applied with \texttt{Fintype.card S}, and
\texttt{if 0 < β then … else …} selecting the sign; \texttt{vecVar phat Zt} is $\Var_{\phat}(\Zt)$.
The function is total in Lean and takes the value $0$ at \texttt{n = 0} (Lean's $\alpha(0)/0 = 0$);
this value never appears in a statement, since every statement about the bonus assumes
\texttt{0 < n} (deviation~1). The constants $2\sqrt2$, $5$, $4H$ are those of Lemma~7 of the
paper (Lemma~\ref{lem:concentration_pos}), of which the bonus is by definition the right-hand
side minus the last term; they are verified there.

\begin{definition}[Optimistic and pessimistic backups]\label{def:backups}
\lean{Essakine2026Tight.stepU, Essakine2026Tight.optZStep, Essakine2026Tight.optZ, Essakine2026Tight.stepUAt, Essakine2026Tight.ZtildeAt, Essakine2026Tight.ZlowerAt, Essakine2026Tight.UtildeAt, Essakine2026Tight.UlowerAt}\leanok
\uses{def:bonus, def:counts, def:emp_trans, def:reward_fn}
Let $t \ge 0$. The \emph{optimistic and pessimistic exponential values} $\Zt^t_h, \Zl^t_h : \Ss \to \R$
($h \le H + 1$) and \emph{backups} $\Ut^t_h, \Ul^t_h : \Ss \times \As \to \R$ ($h \le H$) are defined by the
backward recursion $\Zt^t_{H+1} := 1$, $\Zl^t_{H+1} := 1$ and, for $h = H, \dots, 1$, with the
\emph{clip} $c_h := e^{\beta(H + 1 - h)}$:

\emph{Case $\beta > 0$} (eq.~(optimism\_equation\_positive)). For a pair with $n = n_h^t(s, a) \ge 1$
and $b = b_h^t(s, a)$,
\begin{align*}
  \Ut^t_h(s, a) &:= \min\Big\{ c_h,\ e^{\beta r_h(s, a)} \Big[ \phat \Zt^t_{h+1} + b + \frac1H\, \phat\big(\Zt^t_{h+1} - \Zl^t_{h+1}\big) \Big] \Big\}, \\
  \Ul^t_h(s, a) &:= \max\Big\{ 1,\ e^{\beta r_h(s, a)} \Big[ \phat \Zl^t_{h+1} - b - \frac1H\, \phat\big(\Zt^t_{h+1} - \Zl^t_{h+1}\big) \Big] \Big\},
\end{align*}
and for a never-visited pair ($n_h^t(s, a) = 0$) $\Ut^t_h(s, a) := c_h$, $\Ul^t_h(s, a) := 1$; then
$\Zt^t_h(s) := \max_{a} \Ut^t_h(s, a)$ and $\Zl^t_h(s) := \max_a \Ul^t_h(s, a)$.

\emph{Case $\beta < 0$} (eq.~(optimism\_equation\_negative)). The clips are exchanged: for $n \ge 1$,
\begin{align*}
  \Ut^t_h(s, a) &:= \min\Big\{ 1,\ e^{\beta r_h(s, a)} \Big[ \phat \Zt^t_{h+1} + b + \frac1H\, \phat\big(\Zt^t_{h+1} - \Zl^t_{h+1}\big) \Big] \Big\}, \\
  \Ul^t_h(s, a) &:= \max\Big\{ c_h,\ e^{\beta r_h(s, a)} \Big[ \phat \Zl^t_{h+1} - b - \frac1H\, \phat\big(\Zt^t_{h+1} - \Zl^t_{h+1}\big) \Big] \Big\},
\end{align*}
for a never-visited pair $\Ut^t_h(s, a) := 1$, $\Ul^t_h(s, a) := c_h$; and
$\Zt^t_h(s) := \min_a \Ut^t_h(s, a)$, $\Zl^t_h(s) := \min_a \Ul^t_h(s, a)$.
\end{definition}

\begin{remark}[Deviations 1 and 3: unvisited pairs and the clips]\label{rem:algorithm_clips}
(a) The paper handles never-visited pairs through the convention $\alpha(0)/0 = \infty$, which
makes the bonus infinite and the minimum/maximum equal to the clips; the blueprint defines the
backups of such pairs directly as the clips, so that Lean's $\alpha(0)/0 = 0$ never enters.
(b) The paper writes the optimistic clip as $e^{\beta(H - h - 1)}$ (both in Section~4 and in
eq.~(optimism\_equation\_positive)) and, for $\beta < 0$, the pessimistic clip as
$e^{\beta(H - h - 1)}$. With the paper's $1$-indexed steps this is off by two: the target of the
optimistic backup is $U^*_h$, which ranges in $[1, e^{\beta(H + 1 - h)}]$ for $\beta > 0$
(Lemma~\ref{lem:opt_exp_value_range}), and optimism ($U^*_h \le \Ut^t_h$, Lemma~\ref{lem:optimism_pos})
forces the clip to dominate $e^{\beta(H + 1 - h)}$: at the last step $h = H$ the paper's clip
$e^{-\beta} < 1 \le U^*_H(s, a) = e^{\beta r_H(s, a)}$ would break optimism for every pair. The
recursion is consistent with the clip $c_h = e^{\beta(H + 1 - h)}$: the bracket of $\Ut^t_h$ is
at most $e^{\beta}\cdot e^{\beta(H - h)}$ plus the bonus terms, and $e^{\beta(H + 1 - h)}$ is the
exact range of the values it approximates (Lemma~\ref{lem:backups_range}). In Lean the clip is
\texttt{Real.exp (β * (k + 1))} for \texttt{k = H - h} steps after the step (the paper's
exponent $H - h - 1$ is $k - 1$).
\end{remark}

\textbf{Lean remark.} \texttt{stepU nA H β δ k n r phat Zt Zl : ℝ × ℝ} computes the pair
$(\Ut, \Ul)(s, a)$ of one pair from its step data, with the clips
\texttt{up := if 0 < β then Real.exp (β * (k + 1)) else 1},
\texttt{low := if 0 < β then 1 else Real.exp (β * (k + 1))} and \texttt{if n = 0 then (up, low) else (min up …, max low …)};
\texttt{optZ r β δ hist : ℕ → (S → ℝ) × (S → ℝ)} is the backward recursion
($\Zt, \Zl$ at the step with \texttt{j} steps to the end; \texttt{optZ 0 = (1, 1)}, and
\texttt{optZ (k + 1)} takes the \texttt{max} (\texttt{0 < β}) or \texttt{min} over
\texttt{A} of the \texttt{stepU} of the step \texttt{stepOf H k}); \texttt{stepUAt r β δ hist h s a}
is the pair of backups at the step \texttt{h : Fin H}. The finite maxima and minima are
\texttt{Finset.univ.sup'}/\texttt{inf'} or LML's \texttt{Function.max}/\texttt{min} on a
finite nonempty type.

\begin{lemma}[Ranges and order of the backups]\label{lem:backups_range}
\uses{def:backups, def:bonus, lem:rates_props, lem:emp_trans_simplex}
For every $t$, $h \le H$, $(s, a)$ with $n_h^t(s, a) \ge 1$: $b_h^t(s, a) \ge 0$. For every $t$,
$h \le H$ and $(s, a)$:
\begin{enumerate}
  \item[(i)] if $\beta > 0$: $1 \le \Ul^t_h(s, a) \le \Ut^t_h(s, a) \le e^{\beta(H + 1 - h)}$;
  \item[(ii)] if $\beta < 0$: $e^{\beta(H + 1 - h)} \le \Ul^t_h(s, a) \le \Ut^t_h(s, a) \le 1$;
\end{enumerate}
and for every $h \le H + 1$ and $s$, the same bounds hold for $\Zl^t_h(s) \le \Zt^t_h(s)$:
$1 \le \Zl^t_h(s) \le \Zt^t_h(s) \le e^{\beta(H + 1 - h)}$ if $\beta > 0$, and
$e^{\beta(H + 1 - h)} \le \Zl^t_h(s) \le \Zt^t_h(s) \le 1$ if $\beta < 0$ (at $h = H + 1$ all
values are $1$). In particular $\Zt^t_{h+1}$ and $\Zl^t_{h+1}$ take values in $[1, e^{\beta(H - h)}]$
($\beta > 0$), resp.\ $[e^{\beta(H - h)}, 1]$ ($\beta < 0$), and $\Zt^t_h, \Zl^t_h > 0$.
\end{lemma}

\begin{proof}
\uses{def:backups, def:bonus, lem:rates_props, lem:emp_trans_simplex}
Fix $t$. We prove by backward induction on $h = H + 1, H, \dots, 1$ the statement $P(h)$:
\emph{for all $s$, $\Zl^t_h(s) \le \Zt^t_h(s)$ and both lie in $I_h := [\min(1, e^{\beta(H + 1 - h)}), \max(1, e^{\beta(H + 1 - h)})]$},
proving along the way the claims about $b$, $\Ul$ and $\Ut$ at the step $h$. $P(H + 1)$: both
values are $1 \in I_{H+1} = \{1\}$. Assume $P(h + 1)$ for some $h \le H$ and fix $(s, a)$.

\emph{The bonus is nonnegative.} For $n \ge 1$: $\Var_{\phat}(\cdot) \ge 0$ (Lemma~\ref{lem:emp_trans_simplex}),
$\alpha(n), \alpha^*(n) \ge 0$ (Lemma~\ref{lem:rates_props}), $e^{\beta(H - h)} > 0$, and for $\beta < 0$
also $1 - e^{\beta(H - h)} \ge 0$; so $b_h^t(s, a)$ is a sum of nonnegative terms.

\emph{The correction term is nonnegative.} $w := \frac1H \phat(\Zt^t_{h+1} - \Zl^t_{h+1}) \ge 0$ by
$P(h + 1)$ and the positivity of $\phat$ (Lemma~\ref{lem:emp_trans_simplex}).

\emph{Case $\beta > 0$.} Put $c := e^{\beta(H + 1 - h)} \ge 1$ (as $H + 1 - h \ge 1$). If $n = 0$,
$\Ul = 1 \le c = \Ut$. If $n \ge 1$, let $X := e^{\beta r_h(s, a)}[\phat\Zt^t_{h+1} + b + w]$ and
$Y := e^{\beta r_h(s, a)}[\phat\Zl^t_{h+1} - b - w]$, so that $\Ut = \min\{c, X\}$, $\Ul = \max\{1, Y\}$.
Then: $X \ge 1$, since $e^{\beta r} \ge 1$, $\phat \Zt^t_{h+1} \ge 1$ (by $P(h + 1)$, $\Zt^t_{h+1} \ge 1$,
and $\phat$ preserves lower bounds) and $b, w \ge 0$; $Y \le c$, since
$Y \le e^{\beta r_h(s, a)} \phat \Zl^t_{h+1} \le e^{\beta} e^{\beta(H - h)} = c$ (as $r \le 1$ and
$\Zl^t_{h+1} \le \Zt^t_{h+1} \le e^{\beta(H - h)}$); and $Y \le X$ (as $\Zl^t_{h+1} \le \Zt^t_{h+1}$ and
$b, w \ge 0$). Hence $\max\{1, Y\} \le \min\{c, X\}$: indeed $1 \le c$, $1 \le X$, $Y \le c$ and
$Y \le X$. So $1 \le \Ul \le \Ut \le c$. Taking the maximum over $a$ of both sides preserves
the order and the bounds: $1 \le \Zl^t_h(s) = \max_a \Ul \le \max_a \Ut = \Zt^t_h(s) \le c$, which is $P(h)$.

\emph{Case $\beta < 0$.} Put $c := e^{\beta(H + 1 - h)} \le 1$. If $n = 0$, $\Ul = c \le 1 = \Ut$. If
$n \ge 1$, with $X, Y$ as above, $\Ut = \min\{1, X\}$ and $\Ul = \max\{c, Y\}$. Then $X \ge c$: since
$\beta < 0$ and $r \le 1$, $e^{\beta r} \ge e^{\beta}$, and $\phat\Zt^t_{h+1} \ge e^{\beta(H - h)}$ by
$P(h + 1)$, so $X \ge e^{\beta} e^{\beta(H - h)} + 0 = c$; $Y \le 1$: $e^{\beta r} \le 1$ and
$\phat\Zl^t_{h+1} \le 1$; $Y \le X$ as before; and $c \le 1$. Hence $c \le \Ul \le \Ut \le 1$, and the
minimum over $a$ gives $P(h)$.

The final claims restate $P(h + 1)$ and positivity of the lower bounds.
\end{proof}

\textbf{Lean remark.} The induction is on \texttt{j} for \texttt{optZ r β δ hist j}
(\texttt{optZ\_le}, \texttt{optZ\_mem\_Icc}), and the one-step statement is a lemma about
\texttt{stepU} with hypotheses on \texttt{Zt}, \texttt{Zl} (\texttt{stepU\_fst\_le},
\texttt{stepU\_mem\_Icc}); the bonus nonnegativity is \texttt{bonus\_nonneg} (for
\texttt{0 < n}, or unconditionally, the value at \texttt{n = 0} being $0$).

\section{Greedy policy and certificate}
\label{sec:algorithm_certificate}

\begin{definition}[Greedy policy]\label{def:greedy}
\lean{Essakine2026Tight.greedy, Essakine2026Tight.greedyAt}\leanok
\uses{def:backups}
The policy played at the episode $t + 1$ is the \emph{greedy policy} of the history $\hist_t$:
\[
  \pi^{t+1}_h(s) := \argmax_{a \in \As} \Ut^t_h(s, a) \quad (\beta > 0),
  \qquad
  \pi^{t+1}_h(s) := \argmin_{a \in \As} \Ul^t_h(s, a) \quad (\beta < 0),
\]
for $h \le H$ and $s \in \Ss$, with a fixed tie-breaking rule. Consequently
$\Zt^t_h(s) = \Ut^t_h(s, \pi^{t+1}_h(s))$ for $\beta > 0$ and $\Zl^t_h(s) = \Ul^t_h(s, \pi^{t+1}_h(s))$ for
$\beta < 0$ (Definition~\ref{def:backups}).
\end{definition}

\begin{remark}[Deviation 2: the greedy policy for $\beta < 0$]\label{rem:algorithm_greedy_neg}
The algorithm box of the paper (Section~4) writes $\argmin_a \Ut^t_h(s, a)$ for $\beta < 0$, while
its Appendix~B.2 (after eq.~(optimism\_equation\_negative)) defines $\pi^{t+1}_h(s) = \argmin_a \Ul^t_h(s, a)$
and its analysis uses the latter: the stopping rule compares the certificate to $\Zl^t_1(s_1)$
and Lemma~16 (Lemma~\ref{lem:certificate_neg}) bounds $\Ur - \Ul$ at the greedy action, both
of which need $\Zl^t_h(s) = \Ul^t_h(s, \pi^{t+1}_h(s))$. For $\beta < 0$ the pessimistic $\Ul$ is
the \emph{optimistic} bound on the entropic value $Q^*$ (the exponential is decreasing), so
$\argmin_a \Ul$ is the optimistic greedy choice; this is the definition adopted.
\end{remark}

\textbf{Lean remark.} \texttt{greedy r β δ hist : Policy S A H := fun h s ↦ if 0 < β then argmax (fun a ↦ (stepUAt r β δ hist h s a).1) else argmin (fun a ↦ (stepUAt r β δ hist h s a).2)}
with LML's \texttt{argmax}/\texttt{argmin} on a finite nonempty type (a definite choice of a
maximizer, \texttt{le\_argmax}); \texttt{greedyAt r β δ X Y t ω := greedy r β δ (histAt X Y t ω)}.
The map \texttt{hist ↦ greedy r β δ hist} is measurable because the history type is countable
(\texttt{measurable\_of\_countable}).

\begin{definition}[Certificate]\label{def:certificate}
\lean{Essakine2026Tight.certStep, Essakine2026Tight.cert, Essakine2026Tight.certAt}\leanok
\uses{def:backups, def:bonus, def:greedy}
Let $t \ge 0$ and $\pi := \pi^{t+1}$. The \emph{certificate} of $\pi$ is the family of functions
$(\pi_h \cert^t_h) : \Ss \to \R$, $h \le H + 1$, defined by the backward recursion
$\cert^t_{H+1} := 0$ and, for $h = H, \dots, 1$ and $s \in \Ss$, with $a := \pi_h(s)$,
$n := n_h^t(s, a)$, $\phat := \phat_h^t(\cdot \mid s, a)$ and the clip
\[
  c_h := e^{\beta(H + 1 - h)} \ (\beta > 0), \qquad c_h := 1 \ (\beta < 0):
\]
\[
  (\pi_h \cert^t_h)(s) := \cert^t_h(s, a),
  \qquad
  \cert^t_h(s, a) :=
  \begin{cases}
    \min\Big\{ c_h,\ e^{\beta r_h(s, a)} \Big[ 3\, b_h^t(s, a) + \Big(1 + \frac3H\Big)\, \phat\,\big(\pi_{h+1}\cert^t_{h+1}\big) \Big] \Big\} & \text{if } n \ge 1, \\[4pt]
    c_h & \text{if } n = 0
  \end{cases}
\]
(eq.~(width\_certificate\_positive) and (width\_certificate\_negative) of the paper, with
$\pi_{H+1}\cert^t_{H+1} := 0$). Only the composed functions $\pi_h \cert^t_h$ enter the
recursion and the algorithm; $\cert^t_h(s, a)$ for $a \ne \pi_h(s)$ is never used.
\end{definition}

\begin{remark}[The clip of the certificate]\label{rem:algorithm_cert_clip}
For $\beta > 0$ the paper clips the certificate at $e^{\beta(H - h)}$ in Appendix~B (and at
$e^{\beta(H - h - 1)}$ in Section~4). The certificate must dominate $\Zt^t_h - \Zr^t_h$
(Lemma~\ref{lem:cert_dominates_gap_pos}), which is at most $e^{\beta(H + 1 - h)} - 1$ for a visited
pair and equal to $c_h - \min\{e^{\beta r}(p_h \Zr^t_{h+1}), 1\}$-type quantities for an unvisited
one; with the paper's clip the domination fails at $h = H$ as soon as $e^{\beta} - 1 > 1$. The
clip $c_h = e^{\beta(H + 1 - h)}$ (the same as for $\Ut$, deviation~3) dominates every gap of
Lemma~\ref{lem:backups_range}. For $\beta < 0$ the clip $1$ dominates $1 - e^{\beta(H + 1 - h)}$.
\end{remark}

\textbf{Lean remark.} \texttt{cert r β δ hist : ℕ → S → ℝ}, the backward recursion of the
composed certificate ($\pi_h \cert^t_h$ at the step with \texttt{j} steps to the end;
\texttt{cert 0 = 0}), whose step \texttt{k + 1} computes the greedy action
\texttt{a := greedy r β δ hist h s} at \texttt{h := stepOf H k}, the bonus from
\texttt{optZ r β δ hist k}, the clip \texttt{if 0 < β then Real.exp (β * (k + 1)) else 1}, and
\texttt{if n = 0 then clip else min clip (…)}; \texttt{certAt r β δ X Y t ω h}.

\begin{lemma}[Range of the certificate]\label{lem:certificate_range}
\uses{def:certificate, lem:backups_range, lem:emp_trans_simplex}
For every $t$, $h \le H + 1$ and $s$: $0 \le (\pi^{t+1}_h \cert^t_h)(s) \le c_h$, with
$c_{H+1} := 0$; in particular $\pi^{t+1}_{h+1}\cert^t_{h+1}$ takes values in $[0, e^{\beta(H - h)}]$
($\beta > 0$) or $[0, 1]$ ($\beta < 0$) for $h \le H$.
\end{lemma}

\begin{proof}
\uses{def:certificate, lem:backups_range, lem:emp_trans_simplex}
Backward induction on $h$. At $h = H + 1$ the certificate is $0$. Let $h \le H$ and assume
$0 \le \pi_{h+1}\cert^t_{h+1} \le c_{h+1}$; fix $s$, $a = \pi_h(s)$. If $n = 0$ the value is $c_h \in [0, c_h]$
($c_h > 0$). If $n \ge 1$: $b_h^t(s, a) \ge 0$ (Lemma~\ref{lem:backups_range}),
$\phat(\pi_{h+1}\cert^t_{h+1}) \ge 0$ by the induction hypothesis and positivity of $\phat$
(Lemma~\ref{lem:emp_trans_simplex}), and $e^{\beta r} > 0$, so the bracket is nonnegative and
$\min\{c_h, \cdot\} \in [0, c_h]$. For $\beta > 0$, $c_{h+1} = e^{\beta(H - h)}$; for $\beta < 0$, $c_{h+1} = 1$
(and $0$ at $h = H$, contained in $[0, 1]$).
\end{proof}

\section{Stopping rule and the algorithm}
\label{sec:algorithm_alg}

\begin{definition}[Stopping rule]\label{def:stopping_rule}
\lean{Essakine2026Tight.stopCond}\leanok
\uses{def:certificate, def:backups}
The \emph{stopping condition} after $t$ episodes is
\[
  (\pi^{t+1}_1 \cert^t_1)(s_1) \le \frac{e^{\beta\eps} - 1}{e^{\beta\eps}}\, \Zt^t_1(s_1)
  \quad (\beta > 0, \text{ eq.~(computable\_stopping\_condition)}),
\]
\[
  (\pi^{t+1}_1 \cert^t_1)(s_1) \le \big(1 - e^{\beta\eps}\big)\, \Zl^t_1(s_1)
  \quad (\beta < 0, \text{ Algorithm~1 and Appendix~B.2}).
\]
Entropic-BPI stops after $t$ episodes if the condition holds for the history $\hist_t$ (and not
before). Both sides are computable from $\hist_t$ and $r$ by the backward recursions of
Definitions~\ref{def:backups} and~\ref{def:certificate}.
\end{definition}

\textbf{Lean remark.} \texttt{stopCond r β δ ε s₁ hist : Prop := if 0 < β then cert r β δ hist H s₁ ≤ (Real.exp (β * ε) - 1) / Real.exp (β * ε) * (optZ r β δ hist H).1 s₁ else cert r β δ hist H s₁ ≤ (1 - Real.exp (β * ε)) * (optZ r β δ hist H).2 s₁}
(the step $1$ has \texttt{j = H} steps to the end). For $\beta > 0$ the condition is
equivalent to $\pi\cert \le (e^{\beta\eps} - 1)(\Zt^t_1(s_1) - \pi\cert)$, the form in which it is
used (Lemma~\ref{lem:pac_at_stopping_pos}). Appendix~B.2 of the paper writes the $\beta < 0$
condition as $\pi \cert \le (e^{\beta\eps} - 1)\Zl^t_1(s_1)$, whose right-hand side is negative for
$\beta < 0$ (the algorithm would never stop); the algorithm box and the proof of the sample
complexity use $(1 - e^{\beta\eps})\Zl^t_1(s_1)$, which is the intended condition. Note that the
condition fails after $0$ episodes: all pairs are unvisited, so $\pi^1\cert^0_1(s_1) = c_1$ while
$\Zt^0_1(s_1) = e^{\beta H}$ (resp.\ $\Zl^0_1(s_1) = e^{\beta H}$) and the factors
$(e^{\beta\eps} - 1)/e^{\beta\eps}$, $1 - e^{\beta\eps}$ are less than $1$; hence $\tau \ge 1$.

\begin{definition}[Algorithm 1: Entropic-BPI]\label{def:entropic_bpi}
\lean{Essakine2026Tight.entropicBPI}\leanok
\uses{def:bpi_alg, def:greedy, def:stopping_rule, def:backups, def:certificate}
\emph{Entropic-BPI} with parameters $(r, \beta, \delta, \eps, s_1)$ is the BPI algorithm
(Definition~\ref{def:bpi_alg}) whose
\begin{enumerate}
  \item[(i)] \emph{sampling rule} is deterministic: after $t$ episodes, with history $\hist_t$, it
        plays the greedy policy $\pi^{t+1}$ of $\hist_t$ (Definition~\ref{def:greedy}) at the
        episode $t + 1$;
  \item[(ii)] \emph{stopping rule} is the set of histories $\hist_t$ (of any length $t$) satisfying
        the stopping condition of Definition~\ref{def:stopping_rule};
  \item[(iii)] \emph{output rule} is deterministic: it outputs the greedy policy of the history at
        which it stops, i.e.\ $\hat\pi := \pi^{\tau + 1}$ when it stops after $\tau$ episodes.
\end{enumerate}

\medskip\noindent\textbf{Algorithm 1 (Entropic-BPI).}
\begin{enumerate}
  \item[] \textbf{Input:} the reward function $r$, $\beta \ne 0$, $\delta \in (0, 1)$, $\eps > 0$, the
        initial state $s_1$.
  \item[] \textbf{Initialization:} $t := 0$, empty history ($n_h^0 \equiv 0$, $\phat_h^0 \equiv 1/S$).
  \item[] \textbf{Repeat} (after $t$ episodes, from the history $\hist_t$):
  \item[(1)] \emph{Terminal initialization:} $\Zt^t_{H+1} := 1$, $\Zl^t_{H+1} := 1$, $\cert^t_{H+1} := 0$.
  \item[(2)] \emph{Backward pass, for $h = H, H - 1, \dots, 1$:} for all $(s, a)$ with $n_h^t(s, a) \ge 1$
        compute the bonus $b_h^t(s, a)$ (Definition~\ref{def:bonus}); for all $(s, a)$ compute the
        backups $\Ut^t_h(s, a)$, $\Ul^t_h(s, a)$ and then $\Zt^t_h$, $\Zl^t_h$
        (Definition~\ref{def:backups}); for all $s$ set the greedy action $\pi^{t+1}_h(s)$
        (Definition~\ref{def:greedy}); for all $s$ compute the certificate $(\pi^{t+1}_h\cert^t_h)(s)$
        (Definition~\ref{def:certificate}).
  \item[(3)] \emph{Stopping test:} if the stopping condition of Definition~\ref{def:stopping_rule}
        holds, \textbf{stop}: $\tau := t$, \textbf{output} $\hat\pi := \pi^{t+1}$.
  \item[(4)] \emph{Otherwise} play the episode $t + 1$ with the policy $\pi^{t+1}$ from $s_1$,
        observe its state sequence $(s^{t+1}_1, \dots, s^{t+1}_{H+1})$, append it to the history,
        set $t := t + 1$ and go to (1).
\end{enumerate}
\end{definition}

\textbf{Lean remark.} \texttt{entropicBPI r β δ ε s₁ : BPIAlg S A H} with
\texttt{alg := detAlgorithm (fun \_ p ↦ greedy r β δ p.1) (fun \_ ↦ measurable\_of\_countable \_)}
(LML \texttt{detAlgorithm}: the next action is a measurable function of the pair (history,
observation); the observation is \texttt{()}), \texttt{stopSet := \{h : Σ n, Hist Unit (Policy S A H) (Traj S H) n | stopCond r β δ ε s₁ h.2\}}
(measurable since the history type is countable: \texttt{(Set.to\_countable \_).measurableSet}),
and \texttt{output := Kernel.deterministic (fun h ↦ greedy r β δ h.2) (measurable\_of\_countable \_)}.
This is the LML main branch's \texttt{IdentAlg} with a \texttt{stopSet} of histories of
variable length and one \texttt{output} kernel on them (the previous statement-only
formalization used a stopping predicate \texttt{stop n hist} and a family \texttt{output n}).

\begin{lemma}[Runs of Entropic-BPI]\label{lem:entropic_bpi_run}
\uses{def:entropic_bpi, def:bpi_alg, def:greedy, def:stopping_rule, def:episode_env}
Let $(X, Y, \hat\pi)$ be a run of Entropic-BPI in the episode environment of an MDP
$\mathcal M \in \mathcal M_r$ from $s_1$, on $(\Omega, \Prob)$, and let $\tau$ be its stopping time.
Almost surely:
\begin{enumerate}
  \item[(i)] for every $t \ge 0$, $X_t = \pi^{t+1}$, the greedy policy of the history $\hist_t = (X_i, Y_i)_{i < t}$
        of the first $t$ episodes (Definition~\ref{def:greedy});
  \item[(ii)] $\tau = \min\{t \ge 0 : \hist_t \text{ satisfies the stopping condition}\}$: on
        $\{\tau < \infty\}$ the stopping condition holds for $\hist_\tau$ and fails for $\hist_t$
        for every $t < \tau$; on $\{\tau = \infty\}$ it fails for every $t$;
  \item[(iii)] on $\{\tau < \infty\}$, $\hat\pi = \pi^{\tau + 1}$, the greedy policy of $\hist_\tau$.
\end{enumerate}
\end{lemma}

\begin{proof}
\uses{def:entropic_bpi, def:bpi_alg, def:greedy, def:stopping_rule}
(i) $(X, Y)$ is an algorithm-environment sequence for the sampling rule of Entropic-BPI, a
deterministic algorithm whose next action after the history $\hist_t$ is $\pi^{t+1} = \mathrm{greedy}(\hist_t)$.
The defining property of an algorithm-environment sequence is that the conditional law of
$X_t$ given $\hist_t$ (and the trivial observation) is the policy kernel, here the Dirac mass
at $\mathrm{greedy}(\hist_t)$; a random variable whose conditional law given $Z$ is the Dirac mass
at $g(Z)$ equals $g(Z)$ almost surely. A countable intersection over $t$ gives (i) almost
surely for all $t$ simultaneously.

(ii) The stopping time of an identification algorithm is the hitting time of its stopping rule
by the process $t \mapsto \hist_t$ (Definition~\ref{def:bpi_alg}): $\tau$ is the least $t$ with
$\hist_t$ in the stopping rule, i.e.\ satisfying the stopping condition, and $+\infty$ if there is
none. Minimality gives the failure of the condition for all $t < \tau$, and $\tau < \infty$ gives
$\hist_\tau$ in the stopping rule.

(iii) The output has conditional law $\rho(\hist_\tau)$ given the stopped history $\hist_\tau$,
where $\rho$ is the deterministic kernel $\mathrm{hist} \mapsto \delta_{\mathrm{greedy}(\mathrm{hist})}$; as in (i),
$\hat\pi = \mathrm{greedy}(\hist_\tau)$ almost surely. On $\{\tau < \infty\}$ the stopped history is the
history of the first $\tau$ episodes, so $\hat\pi = \mathrm{greedy}(\hist_\tau) = \pi^{\tau + 1}$
(which is also $X_\tau$ by (i)).
\end{proof}

\textbf{Lean remark.} (i) is LML's \texttt{IsDeterministicAlg.action\_ae\_all\_eq} (or
\texttt{action\_detAlgorithm\_ae\_all\_eq}) for \texttt{detAlgorithm}: \texttt{∀ᵐ ω ∂P, ∀ t, X t ω = greedy r β δ (histAt X Y t ω)};
its proof is \texttt{hasCondDistrib\_action} with the deterministic policy kernel and
\texttt{ae\_eq\_of\_hasCondDistrib\_deterministic}. (ii) is the characterization of
\texttt{hittingAfter} (Mathlib: \texttt{hittingAfter\_mem\_set\_of\_ne\_top},
\texttt{hittingAfter\_le\_iff}, and LML's \texttt{IdentAlg.stoppedHist\_mem\_stopSet\_of\_ne\_top});
the value \texttt{τ ω : ℕ∞} is compared with \texttt{t} through \texttt{ENat} coercions. (iii)
is \texttt{IsRun.hasCondDistrib\_output} with \texttt{Kernel.deterministic}, giving
\texttt{out =ᵐ[P] fun ω ↦ greedy r β δ (A.stoppedHist O X Y ω).2}, and
\texttt{stoppedHist ω = ⟨τ ω, histAt X Y (τ ω) ω⟩} on \texttt{τ ω ≠ ⊤} (LML
\texttt{stoppedValue} of \texttt{sigmaHistory}). Statements about the output are only made on
$\{\tau < \infty\}$, which has full probability under the good event
(Lemma~\ref{lem:stopping_time_bound_pos}).

\section{The ring quantities of the analysis}
\label{sec:algorithm_ring}

The analysis compares the computable values to the values of the played policy through an
auxiliary backward recursion, which follows the true exponential Bellman recursion of the
greedy policy but is clipped by the empirical pessimistic (resp.\ optimistic) backup. It uses
the unknown kernels $p_h$ and is never computed by the algorithm.

\begin{definition}[Ring values and backups]\label{def:ring}
\lean{Essakine2026Tight.ringUAux, Essakine2026Tight.ringZStep, Essakine2026Tight.ringZ, Essakine2026Tight.ringU, Essakine2026Tight.ringZAt, Essakine2026Tight.ringUAt}\leanok
\uses{def:backups, def:bonus, def:greedy, def:episodic_mdp, def:reward_fn}
Let $\mathcal M \in \mathcal M_r$ with transition kernels $(p_h)$, $t \ge 0$ and $\pi := \pi^{t+1}$. The
\emph{ring values} $\Zr^t_h : \Ss \to \R$, $h \le H + 1$, and \emph{ring backups}
$\Ur^t_h : \Ss \times \As \to \R$, $h \le H$, are defined by $\Zr^t_{H+1} := 1$ and, for $h = H, \dots, 1$
and every $(s, a)$, with $n := n_h^t(s, a)$, $\phat := \phat_h^t(\cdot \mid s, a)$, $b := b_h^t(s, a)$
and $e^{\beta r} := e^{\beta r_h(s, a)}$:

\emph{Case $\beta > 0$} (Appendix~B.1, ``Stopping rule''). For $n \ge 1$,
\begin{align*}
  \Ur^t_{h, \mathrm{pes}}(s, a) &:= \max\Big\{ 1,\ e^{\beta r} \Big[ \phat \Zr^t_{h+1} - b - \frac1H\, \phat\big(\Zt^t_{h+1} - \Zr^t_{h+1}\big) \Big] \Big\}, \\
  \Ur^t_h(s, a) &:= \min\Big\{ e^{\beta r}\, (p_h \Zr^t_{h+1})(s, a),\ \Ur^t_{h, \mathrm{pes}}(s, a) \Big\},
\end{align*}
and for a never-visited pair $\Ur^t_h(s, a) := \min\{ e^{\beta r}\, (p_h \Zr^t_{h+1})(s, a),\ 1\}$
(the pessimistic backup being $\Ul^t_h(s, a) = 1$ there).

\emph{Case $\beta < 0$} (Appendix~B.2). For $n \ge 1$,
\begin{align*}
  \Ur^t_{h, \mathrm{opt}}(s, a) &:= \min\Big\{ 1,\ e^{\beta r} \Big[ \phat \Zr^t_{h+1} + b + \frac1H\, \phat\big(\Zr^t_{h+1} - \Zl^t_{h+1}\big) \Big] \Big\}, \\
  \Ur^t_h(s, a) &:= \max\Big\{ e^{\beta r}\, (p_h \Zr^t_{h+1})(s, a),\ \Ur^t_{h, \mathrm{opt}}(s, a) \Big\},
\end{align*}
and for a never-visited pair $\Ur^t_h(s, a) := \max\{ e^{\beta r}\, (p_h \Zr^t_{h+1})(s, a),\ 1\}$
(the optimistic backup being $\Ut^t_h(s, a) = 1$ there).

In both cases $\Zr^t_h(s) := \Ur^t_h(s, \pi_h(s))$. Here $\Zt^t_{h+1}$, $\Zl^t_{h+1}$ and $b$ are those
of Definitions~\ref{def:backups} and~\ref{def:bonus} for the same history, and
$(p_h \Zr^t_{h+1})(s, a) = \sum_{s'} p_h(s' \mid s, a)\Zr^t_{h+1}(s')$ uses the true kernel.
\end{definition}

\textbf{Lean remark.} \texttt{ringUAux nA H β δ k n r p phat Zr Zt Zl : ℝ} is the one-step
function (arguments: the true transition vector \texttt{p = M.transVec h s a}, the empirical
one, the next-step ring value and the next-step optimistic and pessimistic values), with the
sign selected by \texttt{if 0 < β} and the unvisited case by \texttt{if n = 0};
\texttt{ringZ M β δ hist : ℕ → S → ℝ} is the backward recursion ($\Zr$ at the step with
\texttt{j} steps to the end, \texttt{ringZ 0 = 1}, evaluated at the greedy action
\texttt{greedy r β δ hist h s} of the same history, where \texttt{r} is the reward function of
\texttt{M}), \texttt{ringU M β δ hist h s a} the ring backup at the step \texttt{h}, and
\texttt{ringZAt}, \texttt{ringUAt} their values on a run. The lemmas about these quantities
(Lemmas~\ref{lem:ring_pos}, \ref{lem:certificate_pos}, \ref{lem:cert_dominates_gap_pos} and
their negative counterparts) are in Chapters~\ref{chap:analysis_pos} and~\ref{chap:analysis_neg}.
